\subsection{台形公式}

定積分を近似する最初の方法は,\sectionref{節:リーマン和}で見たように,区間を細かく分けて長方形の面積を足すことであった.長方形の高さには小区間の端の値を使うので,曲線が傾いている所では,長方形と曲線の間に三角形に近い隙間が残る.そこで,小区間の両端の値を結んで台形を作ってみよう.

区間$[a,b]$で,両端の点を通る直線は
\[
 L(x)=\frac{b-x}{b-a}f(a)+\frac{x-a}{b-a}f(b)
\]
である.この直線の下の符号付き面積は台形の面積であり
\[
 T=\int_a^bL(x)\dd x=\frac{b-a}{2}\{f(a)+f(b)\}
\]
となる.これが1つの区間に対する台形公式である.区間全体を$N$等分し,\,$h=\frac{b-a}{N}$,\,$x_j=a+jh$とすると,各小区間の台形を足した値は
\[
 T_N=h\left\{\frac{f(a)+f(b)}2+\sum_{j=1}^{N-1}f(x_j)\right\}
\]
となる.内部の点は隣り合う2つの台形で半分ずつ使われるため,係数が$1$になる.

ここまでは高校の範囲の計算である.では,長方形を台形に替えると,近似はどれだけ良くなるのだろうか.実験してみよう.\,$\int_0^1\e^x\dd x=\e-1$を,小区間の左端の値を高さとする長方形の和と台形公式$T_N$で近似し,誤差を比べると次の表のようになる.
\begin{center}
\begin{tabular}{ccc}
\toprule
$N$ & 長方形の和の誤差 & 台形公式の誤差$T_N-(\e-1)$\\
\midrule
$1$ & $-0.718282$ & $0.140859$\\
$2$ & $-0.393921$ & $0.035649$\\
$4$ & $-0.205845$ & $0.008940$\\
$8$ & $-0.105156$ & $0.002237$\\
$16$ & $-0.053137$ & $0.000559$\\
\bottomrule
\end{tabular}
\end{center}
ここで,どちらの誤差も「近似値から真の値を引いたもの」である.\,$N$を2倍にすると,長方形の和の誤差はおよそ半分になるが,台形公式の誤差はおよそ$\frac14$になる.台形公式の誤差は$N^2$にほぼ反比例すると予想できるだろう.また,台形公式の値はどれも真の値より大きい.

大小関係の方は図から分かる.図\ref{図:台形近似}では,曲線が下に凸なので弦が曲線より上にある.この場合,台形公式の値は積分の値より大きくなる.ただし,関数によって凸凹が変わる場合もあり,誤差の大きさは図からは分からない.そこで,第2次導関数を使って誤差を式で表すことを目標とする.以下,台形公式の誤差を調べる間は,\,$f$は$[a,b]$で2回連続微分可能であるとする.

\begin{figure}[htbp]
\centering
\begin{tikzpicture}[x=1.7cm,y=1.4cm]
 \fill[blue!10] (0.5,0)--(2.5,0)--(2.5,1.9625)--(0.5,0.4625)--cycle;
 \draw[-{Latex}] (-0.1,0)--(2.9,0) node[right] {$x$};
 \draw[-{Latex}] (0,-0.1)--(0,2.35) node[above] {$y$};
 \draw[dashed] (0.5,0)--(0.5,0.4625);
 \draw[dashed] (2.5,0)--(2.5,1.9625);
 \draw[thick,orange!80!black] (0.5,0.4625)--(2.5,1.9625);
 \draw[thick,blue!70!black,domain=0.15:2.65,samples=80] plot (\x,{0.4+0.25*\x*\x});
 \fill (0.5,0.4625) circle (1.5pt);
 \fill (2.5,1.9625) circle (1.5pt);
 \node[below] at (0.5,0) {$a$};
 \node[below] at (2.5,0) {$b$};
 \node[above left] at (1.55,1.25) {$L(x)$};
 \node[right] at (2.05,1.18) {$f(x)$};
\end{tikzpicture}
\caption{端点を結ぶ直線の積分が台形公式になる.下に凸の曲線では,弦と曲線の間の面積だけ過大になる.}
\label{図:台形近似}
\end{figure}

1つの区間で考える.誤差を第2次導関数で表したいので,\,$f''$を含む積分から$f$の積分を取り出すことを考えよう.そのために,端点で$0$となる二次式
\[
 w(x)=(x-a)(b-x)
\]
を$f''$に掛けて,2回部分積分する.\,$w(a)=w(b)=0$なので1回目の境界の項は消え,\,$w'(x)=a+b-2x$,\,$w''(x)=-2$である.
\begin{align*}
 \int_a^b w(x)f''(x)\dd x
 &=-\int_a^bw'(x)f'(x)\dd x\\
 &=-[w'(x)f(x)]_a^b+\int_a^bw''(x)f(x)\dd x\\
 &=(b-a)\{f(a)+f(b)\}-2\int_a^bf(x)\dd x.
\end{align*}
右辺の$(b-a)\{f(a)+f(b)\}$は$2T$である.したがって,誤差を「積分の値から近似値を引いたもの」と決めれば
\[
 \boxed{\int_a^bf(x)\dd x-T
 =-\frac12\int_a^b(x-a)(b-x)f''(x)\dd x}
\]
を得る.\,$a\le x\le b$では$(x-a)(b-x)\ge0$なので,\,$f''\ge0$なら右辺は$0$以下になる.図で見た大小関係が式からも分かるのだ.また,区間で$|f''(x)|\le M_2$なら
\[
 \left|\int_a^bf(x)\dd x-T\right|
 \le\frac{M_2}{2}\int_a^b(x-a)(b-x)\dd x
 =\frac{M_2(b-a)^3}{12}.
\]

$N$等分した場合は,各小区間の誤差が$\frac{M_2h^3}{12}$以下なので,合計では
\[
 \left|\int_a^bf(x)\dd x-T_N\right|
 \le\frac{M_2(b-a)h^2}{12}
 =\frac{M_2(b-a)^3}{12N^2}
\]
となる.分割数を2倍にすれば,この誤差の上限は$\frac14$になる.\,$\e^x$の例では$M_2=\e$とでき,\,$N=4$での上限はおよそ$0.0142$である.表の誤差$0.008940$はこの範囲に収まっている.

\paragraphbf{誤差が$N^2$に反比例すること}
上の評価は誤差の上限であり,誤差そのものが$N^2$に反比例することまでは言っていない.表の予想を証明するため,もう少し詳しく調べよう.\,$j$番目の小区間$[x_{j-1},x_j]$での$f''$の最小値を$m_j$,最大値を$M_j$とする.\,$(x-x_{j-1})(x_j-x)\ge0$と$\int_{x_{j-1}}^{x_j}(x-x_{j-1})(x_j-x)\dd x=\frac{h^3}{6}$から,この小区間の誤差$E_j$は
\[
 -\frac{h^3}{12}M_j\le E_j\le-\frac{h^3}{12}m_j
\]
を満たす.\,$j=1,\ldots,N$について足し,\,$N^2h^2=(b-a)^2$を使うと
\[
 -\frac{(b-a)^2}{12}\sum_{j=1}^{N}M_jh
 \le N^2\left(\int_a^bf(x)\dd x-T_N\right)
 \le-\frac{(b-a)^2}{12}\sum_{j=1}^{N}m_jh
\]
となる.両端の和は,最小値・最大値をとる点を代表点としたリーマン和なので,\,$N\to\infty$でどちらも$\int_a^bf''(x)\dd x=f'(b)-f'(a)$に近づく.したがって
\[
 N^2\left(\int_a^bf(x)\dd x-T_N\right)
 \longrightarrow-\frac{(b-a)^2}{12}\{f'(b)-f'(a)\}
\]
である.\,$\int_0^1\e^x\dd x$では右辺の符号を変えた$N^2\{T_N-(\e-1)\}$が$\frac{\e-1}{12}=0.1431\cdots$に近づく.実際,表の$N=16$では$16^2\times0.000559=0.143\cdots$である.誤差がおよそ$N^2$に反比例するという予想が証明されたのだ.
